\newpage
\subsubsection{Описание состояний конечного автомата управления}
\label{sec:states}

В каждом состоянии определены показания индикатора и действия кнопок.

\begin{enumerate}[itemsep=0.45\baselineskip]
	\item $S_0$~--- отображение текущего времени. Это начальное состояние. На индикаторе видны часы, минуты и секунды в 24-часовом формате.
	\begin{itemize}
		\item $a$~--- перейти к корректировке секунд ($S_1$).
		\item $b$~--- показать секундомер ($S_4$).
		\item $c$~--- показать время будильника и признак его включения ($S_7$).
		\item $d$~--- сохранить текущее состояние.
	\end{itemize}

	\item $S_1$~--- корректировка секунд. На индикаторе отображается поле секунд.
	\begin{itemize}
		\item $b$~--- увеличить секунды на единицу по модулю $60$, оставаясь в $S_1$.
		\item $a$, $c$, $d$~--- перейти к корректировке минут ($S_2$).
	\end{itemize}

	\item $S_2$~--- корректировка минут. На индикаторе отображается поле минут.
	\begin{itemize}
		\item $b$~--- увеличить минуты на единицу по модулю $60$, оставаясь в $S_2$.
		\item $a$, $c$, $d$~--- перейти к корректировке часов ($S_3$).
	\end{itemize}

	\item $S_3$~--- корректировка часов. На индикаторе отображается поле часов.
	\begin{itemize}
		\item $b$~--- увеличить часы на единицу по модулю $24$, оставаясь в $S_3$.
		\item $a$, $c$, $d$~--- завершить корректировку и остановить ход часов ($S_{10}$).
	\end{itemize}

	\item $S_4$~--- отображение секундомера. На индикаторе видны минуты и секунды отсчета. Признак хода секундомера сохраняется при входе в этот режим из $S_0$.
	\begin{itemize}
		\item $a$~--- запустить или продолжить отсчет с сохраненного значения ($S_5$).
		\item $b$~--- остановить секундомер и зафиксировать показание ($S_6$).
		\item $c$, $d$~--- вернуться к текущему времени ($S_0$), сохранив признак хода секундомера.
	\end{itemize}

	\item $S_5$~--- запущенный секундомер. На индикаторе обновляется показание секундомера. Любая кнопка возвращает автомат в $S_4$ при продолжающемся отсчете.

	\item $S_6$~--- остановленный секундомер. На индикаторе остается зафиксированное показание.
	\begin{itemize}
		\item $c$~--- сбросить показание до $00{:}00$, оставаясь в $S_6$.
		\item $a$, $b$, $d$~--- вернуться в $S_4$ при остановленном отсчете.
	\end{itemize}

	\item $S_7$~--- настройка будильника. На индикаторе видны установленное время будильника и признак его включения.
	\begin{itemize}
		\item $a$~--- переключить признак включения будильника, оставаясь в $S_7$.
		\item $b$~--- перейти к корректировке минут будильника ($S_8$).
		\item $c$~--- показать номер мелодии будильника ($S_{11}$).
		\item $d$~--- вернуться к отображению текущего времени ($S_0$).
	\end{itemize}

	\item $S_8$~--- корректировка минут будильника. На индикаторе отображаются минуты установленного времени.
	\begin{itemize}
		\item $a$~--- увеличить минуты на единицу по модулю $60$, оставаясь в $S_8$.
		\item $b$, $c$, $d$~--- перейти к корректировке часов будильника ($S_9$).
	\end{itemize}

	\item $S_9$~--- корректировка часов будильника. На индикаторе отображаются часы установленного времени.
	\begin{itemize}
		\item $a$~--- увеличить часы на единицу по модулю $24$, оставаясь в $S_9$.
		\item $b$, $c$, $d$~--- вернуться к настройкам будильника ($S_7$).
	\end{itemize}

	\item $S_{10}$~--- останов часов. Индикатор показывает время на момент останова. Нажатие любой кнопки возобновляет ход часов и возвращает автомат в $S_0$.

	\item $S_{11}$~--- выбор мелодии будильника. На индикаторе отображается номер выбранной мелодии.
	\begin{itemize}
		\item $a$~--- выбрать следующую мелодию, оставаясь в $S_{11}$. После последней мелодии выбирается первая.
		\item $b$, $c$, $d$~--- вернуться к настройкам будильника ($S_7$).
	\end{itemize}
\end{enumerate}

Вершина $S_{12}$ исходного рисунка соответствует активному состоянию звукового блока. В начале очередной четверти часа автомат выдает команду $z_{12}$, и блок включает сигнал на одну секунду. По истечении секунды блок возвращается в режим ожидания. Состояние автомата управления и назначение кнопок при этом сохраняются.

Коды состояний автомата управления приведены в таблице~\ref{tab:states}. Для 12 состояний необходимо $\lceil\log_2 12\rceil = 4$ разряда. Код $q_3q_2q_1q_0$ совпадает с двоичной записью номера состояния, $q_3$~--- старший разряд.

\begin{table}[H]
	\centering
	\caption{Состояния и их двоичные коды}
	\label{tab:states}
	\small\setstretch{1.12}
	\renewcommand{\arraystretch}{1.12}
	\begin{tabularx}{\textwidth}{|c|>{\raggedright\arraybackslash}p{6.2cm}|>{\raggedright\arraybackslash}X|c|}
		\hline
		Состояние & Режим & Показание индикатора & Код \\
		\hline
		$S_0$ & Отображение времени & Часы, минуты и секунды & $0000$ \\
		$S_1$ & Корректировка секунд & Секунды & $0001$ \\
		$S_2$ & Корректировка минут & Минуты & $0010$ \\
		$S_3$ & Корректировка часов & Часы & $0011$ \\
		$S_4$ & Отображение секундомера & Минуты и секунды отсчета & $0100$ \\
		$S_5$ & Запущенный секундомер & Текущее показание & $0101$ \\
		$S_6$ & Остановленный секундомер & Зафиксированное показание & $0110$ \\
		$S_7$ & Настройка будильника & Время и признак включения & $0111$ \\
		$S_8$ & Установка минут будильника & Минуты будильника & $1000$ \\
		$S_9$ & Установка часов будильника & Часы будильника & $1001$ \\
		$S_{10}$ & Остановленные часы & Время на момент останова & $1010$ \\
		$S_{11}$ & Выбор мелодии будильника & Номер мелодии & $1011$ \\
		\hline
	\end{tabularx}
\end{table}

